\section{The Registers, and How Few There Are}
\label{sec:gpio-registers}

A GPIO port on the EFR32MG24 is not a black box. It is five things:

\begin{table}[htbp]
    \centering
    \footnotesize
    \begin{tabular}{@{}p{0.18\textwidth}p{0.42\textwidth}p{0.30\textwidth}@{}}
        \toprule
        \textbf{Register} & \textbf{What it decides} & \textbf{Access} \\
        \midrule
        \codeinline{MODEL} & mode of pins 0--7 & read-modify-write, per pin \\
        \codeinline{MODEH} & mode of pins 8--15 & only on PORTA and PORTC here \\
        \codeinline{DOUT} & the output latch, one bit per pin & read-modify-write \\
        \codeinline{DIN} & the pin as the outside world holds it & read only \\
        \codeinline{CMU\_CLKEN0} & whether the port answers at all & gate before the first access \\
        \bottomrule
    \end{tabular}
    \caption{The GPIO port, in full. Five registers, and only one of them is optional.}
    \label{tab:gpio-registers}
\end{table}

The modes are numbers, not types:

\begin{minted}{rust}
// Measured from the PAC's field encodings, and confirmed by the effect on the pin.
const DISABLED: u32 = 0;   // the pad is off; an analog peripheral may claim it
const INPUT:    u32 = 1;   // input buffer on, no pull
const PUSH_PULL:u32 = 4;   // driven high or low
const WIRED_AND: u32 = 8;  // open drain: 0 pulls low, 1 lets the pull-up win
const WIRED_OR:  u32 = 9;  // open drain, inverted sense
\end{minted}

Two of these decide most of the code in this book. \codeinline{PUSH_PULL} drives an LED and a
piezo. \codeinline{WIRED_AND} is what makes a bit-banged I2C bus possible, because ``output 1''
means ``release the line'' and the external pull-up makes it a real 1 — which is exactly what an
open-drain bus requires and exactly what a push-pull pin must never do.

\subsection{The Clock Gate Comes First}

The most expensive mistake in this whole project was not a wrong bit. It was touching a register
whose bus clock was off:

\begin{minted}{rust}
// Before: the write faulted, and the fault handler did not exist yet.
gpio.porta_dout().modify(|_, w| w.dout7().set_bit());

// After: gate the port, then talk to it.
dp.cmu_s.clken0().modify(|_r, w| w.gpio().set_bit());
gpio.porta_dout().modify(|_, w| w.dout7().set_bit());
\end{minted}

The measurement, taken from the fault status register after the fact:
\codeinline{CFSR = 0x8200} (a precise bus fault) with \codeinline{BFAR} pointing at the GPIO address
that was written. A GPIO register with no clock is not ``reads as zero'' — it is a bus fault, and if
the vector table still points at a default handler the board simply stops.

\subsection{Not Every Port Has Every Register}

\begin{minted}{rust}
// This package is a 48-pin QFN: PORTB and PORTD expose eight pins, not sixteen.
// MODEH covers 8..15, so on those two ports it does not exist at all.
match port {
    Port::A | Port::C => { /* MODEL and MODEH */ }
    Port::B | Port::D => { /* MODEL only */ }
}
\end{minted}

That is not a simplification invented for this book; it is what the SVD says, and a HAL that
pretends otherwise fails to compile \emph{after} someone has written the code. Two of the eight
compile errors during this project's pin primitive came from exactly that: the generated register
blocks differ per port, and a macro written for one does not typecheck for the other.
